\section{Introduction}
Functional enrichment analysis is the default first step in interpreting gene lists produced by genetic association studies, differential-expression experiments or curated disease resources. The most widely used approach, over-representation analysis (ORA), asks whether the query list shares more genes with a pathway than expected under a hypergeometric model \citep{Huang2009}. ORA has three well-documented limitations. First, it only uses \emph{shared} genes: a query gene that is not annotated to a pathway contributes nothing, even if it interacts directly with the pathway's core members. Pathway annotations are, however, incomplete and strongly biased towards well-studied genes \citep{Haynes2018}, and genes implicated by genome-wide association studies (GWAS) frequently act on pathways through interaction partners rather than as annotated members \citep{Barabasi2011,Cowen2017}. Second, the hypergeometric null assumes that every gene is equally likely to be selected and annotated. Highly connected, well-studied proteins are over-represented both among disease genes and among pathway members, which inflates ORA's false-positive rate \citep{Timmons2015,Haynes2018}; an analogous selection bias (gene length in RNA-seq) motivated GOseq \citep{Young2010}. Third, results depend strongly on analysis choices and are often reported without calibration \citep{Wijesooriya2022}.

Network-based enrichment methods use a protein--protein interaction (PPI) network to recover the information that ORA discards. Network enrichment analysis (NEA) and its successors BinoX and ANUBIX count edges between the query and a pathway \citep{Alexeyenko2012,Ogris2017,CastresanaAguirre2020}; EnrichNet and propagation-based methods diffuse the query over the network \citep{Glaab2012,Cowen2017}; pathfindR searches for active subnetworks before running ORA \citep{Ulgen2019}; and topology-based methods such as SPIA use the internal wiring of pathways rather than a genome-wide interactome \citep{Ma2019}. In network medicine, the shortest-path \emph{proximity} of two gene modules, assessed against a degree-preserving random model, has been the workhorse for relating diseases to each other and to drug targets \citep{Menche2015,Guney2016,MorselliGysi2021}. Applying proximity to pathways is natural, but existing implementations, including an earlier, unpublished implementation of our own, share three problems that limit their adoption for enrichment analysis: (i) the proximity score mixes overlap with network evidence, because shared genes have distance zero, so it is unclear whether a network method finds anything that ORA does not; (ii) significance is obtained from a limited number of permutations summarised by a normal approximation, which is slow, has limited resolution in the tail where multiple-testing thresholds lie, and is not guaranteed to be calibrated; and (iii) the interactome itself often contains edges derived from pathway databases (e.g.\ the curated-database channel of STRING), making network proximity to a pathway partly circular.

Here we present NetFunrankR, which addresses all three problems with a single, statistically explicit construction. Our contributions are:
\begin{enumerate}[leftmargin=*,itemsep=2pt]
\item \textbf{A decomposition of network proximity into two independent tests.} The closest-distance proximity of a query to a gene set is split into a degree-aware \emph{overlap} test and \emph{beyond-overlap} tests that evaluate only the non-member query genes---the number of query genes that interact directly with the pathway (linker genes) and their distance to it---conditionally on the membership pattern. We show that the two p-values are independent under the null, so that they can be combined exactly by Fisher's method and each enriched pathway can be attributed to overlap, network proximity or both (Proposition~\ref{prop:indep}).
\item \textbf{Exact, permutation-free degree-preserving null distributions.} Under the standard degree-binned null of network medicine, all three statistics are sums of independent discrete variables whose distributions we compute exactly by convolution (and by a saddlepoint approximation for long lists). This removes the permutation step, gives accurate p-values far into the tail, and reduces run time from hours to seconds.
\item \textbf{A degree-aware overlap test} that corrects the hub bias of the hypergeometric test, and a \textbf{non-circular interactome} (STRING v12 recombined without database and text-mining evidence).
\item \textbf{A comprehensive, multi-scenario validation} against ORA, nine established expression-based methods and four network-based methods: type-I error calibration under realistic hub-biased null lists, power under controlled annotation incompleteness, the GSEABenchmarkeR compendium of 42 expression datasets \citep{Geistlinger2021}, GWAS gene lists for 25 diseases with three independent gold standards including a temporal (prospective) validation, robustness to interactome choice and perturbation, and a COPD case study.
\end{enumerate}
